\documentclass[11pt,a4paper]{article}

% ── Packages ──
\usepackage[margin=2.5cm]{geometry}
\usepackage[utf8]{inputenc}
\usepackage[T1]{fontenc}
\usepackage{amsmath,amssymb}
\usepackage{booktabs}
\usepackage{longtable}
\usepackage{array}
\usepackage{graphicx}
\usepackage[dvipsnames]{xcolor}
\usepackage{tikz}
\usetikzlibrary{arrows.meta, positioning, shapes.geometric, fit, calc}
\usepackage[numbers,sort&compress]{natbib}
\usepackage{hyperref}
\hypersetup{colorlinks=true,linkcolor=blue,citecolor=blue,urlcolor=blue}
\usepackage[protrusion=true,expansion=false]{microtype}
\usepackage{enumitem}
\setlist{nosep}

% ── Layout tweaks ──
\setlength{\parindent}{0pt}
\setlength{\parskip}{0.6em}
\renewcommand{\baselinestretch}{1.15}

% ── Custom commands ──
\newcommand{\HRule}{\rule{\linewidth}{0.4pt}}

% ══════════════════════════════════════════════════════════════════
\begin{document}

% ── Title ──
\begin{center}
{\LARGE\bfseries Hierarchical glycan prediction from protein sequence\\
reveals a determination boundary at the structural-class level}

\vspace{1em}
{\large Yusuke Matsui$^{1,2,*}$}\\[0.3em]
{\normalsize $^1$Biomedical and Health Informatics Unit, Department of Integrated Health Sciences,\\
Nagoya University Graduate School of Medicine, Nagoya, Japan\\
$^2$Institute for Glyco-core Research (iGCORE), Nagoya University, Nagoya, Aichi, Japan}\\[0.3em]
{\normalsize $^*$Corresponding author: Y.\ Matsui}
\end{center}
\HRule
\vspace{0.5em}

% ══════════════════════════════════════════════════════════════════
\section*{Abstract}

\textbf{Motivation:} Predicting which glycan structures are attached at specific
protein sites remains an open challenge. Glycobiology data are fragmented
across specialised databases, and the relative contributions of protein
sequence versus cellular environment to glycan identity have not been
quantified.

\textbf{Results:} We present glycoMusubi, an integrated knowledge graph spanning
78,263 nodes (10 entity types) and 2.5 million edges (14 relation types) from
six public databases. Using glycoMusubi, we establish a hierarchical benchmark
decomposing glycan prediction into six tasks of increasing specificity, from
N-linked site classification (F1 = 0.924) through structural cluster prediction
(Hits@5 = 0.786) to exact glycan identification (MRR = 0.066--0.118). The
10-fold gap between cluster-level and exact-ID performance reveals an
empirical determination boundary: protein sequence predicts structural class
reliably, but exact glycan identity requires additional information. Integration
of single-cell enzyme expression data (Tabula Sapiens; 748 cell types) did not
improve prediction beyond a frequency baseline, because current enzyme--glycan
annotations cover only 39\% of N-linked glycans. These results delineate two
layers of glycosylation determination and identify annotation sparsity as the
primary bottleneck.

\textbf{Availability and implementation:}
The glycoMusubi software and benchmark suite are available at
\url{https://github.com/ymatts/glycoMusubi} under the MIT licence.
The knowledge graph and trained models will be deposited on Zenodo.

\textbf{Contact:} Y.\ Matsui

\textbf{Supplementary information:} Supplementary data are available at
\textit{Bioinformatics} online.

\HRule

% ══════════════════════════════════════════════════════════════════
\section{Introduction}

Glycans---complex carbohydrates covalently attached to proteins and
lipids---mediate virtually every cell-surface biological process, from immune
recognition and cell signalling to pathogen entry and cancer
progression~\cite{varki2017,reily2019}. Despite decades of structural
characterisation, a fundamental question remains unanswered computationally:
given a protein sequence and a cellular context, which glycan structures will
be attached at specific sites?

This question is challenging because glycosylation is determined at multiple
levels. The protein sequence and local site context influence which
glycosyltransferases can access the substrate, constraining the structural
class of possible glycans~\cite{schjoldager2020,narimatsu2019}. However, the
exact glycan structure further depends on the complement of glycosylation
enzymes expressed in the cell, the availability of nucleotide sugar donors and
the kinetics of competing biosynthetic
pathways~\cite{stanley2022,varki2022essentials}. The relative contribution of
these `sequence-determined' versus `environment-determined' factors has been
discussed qualitatively~\cite{schjoldager2020} but never quantified
systematically.

Recent computational approaches have tackled parts of this problem.
InSaNNE~\cite{kellman2024} trained a recurrent neural network to predict glycan
structures from protein sequence features, achieving AUC = 0.92 on a curated set
of 199 glycans across 57 glycosites. However, the approach evaluates binary
presence/absence rather than retrieval ranking, uses a restricted candidate set
and does not decompose the hierarchy of prediction difficulty.
Bojar and Lisacek~\cite{bojar2022} reviewed glycoinformatics and machine learning
approaches, noting the gap between glycan property prediction and the inverse
problem of predicting glycans from protein context.
CoNglyPred~\cite{congly2025} combined ESM-2 and AlphaFold2 features for N-linked
site prediction, complementing our Task~1. GlycanML~\cite{glyml2024} established
a benchmark of 11 glycan property prediction tasks (immunogenicity, taxonomy,
tissue expression), but these predict \textit{from} glycan structures rather than
\textit{predicting} glycan structures from protein context.

On the biological side, Yadav \textit{et al.}~\cite{yadav2022} modelled Golgi
glycan processing as an optimal information-coding channel, showing that
structural diversity is bounded by enzyme kinetics---a theoretical framework
consistent with our empirical determination boundary.
Chrysinas \textit{et al.}~\cite{chrysinas2024} analysed glycosylation enzyme
expression from the same Tabula Sapiens atlas we use, finding that terminal
processing enzymes are highly cell-type-specific while core pathway enzymes are
ubiquitous. Their qualitative observation aligns with our quantitative finding
that enzyme expression alone does not resolve exact glycan identity.
Glycoimpact~\cite{glycoimpact2024} defined an amino acid substitution scoring
matrix for glycan constraints, introducing the concept of `bounded biosynthesis'
where sequence constrains but does not determine the glycan.

Progress has also been hindered by data fragmentation. Glycan structures reside
in GlyTouCan~\cite{tiemeyer2017}, protein annotations in
UniProt~\cite{uniprot2023}, enzyme--glycan biosynthetic relationships in
GlyGen~\cite{york2020}, inhibitor activities in ChEMBL~\cite{zdrazil2024} and
site-specific glycan assignments in GlyConnect~\cite{alocci2019}. Each database
uses distinct schemas and identifiers, making integrated analysis difficult.
Recent advances in glycan representation---SweetNet~\cite{burkholz2022} for graph
neural network encoding, GIFFLAR~\cite{thomes2024} for higher-order
representations and glycowork for motif-based analysis~\cite{thomes2021}---have
improved glycan property prediction but have not addressed the inverse problem of
predicting glycans from protein context, nor quantified the determination
boundary.

Here we make three contributions centred on resources and biological
insight rather than methodological novelty.  First, we construct
\textbf{glycoMusubi}, a heterogeneous knowledge graph integrating 78,263
nodes (10 types) and 2,508,960 edges (14 relation types) from six public
databases through an automated, reproducible ETL pipeline.  Second, we
establish a \textbf{hierarchical prediction benchmark} comprising six
tasks of increasing specificity, from binary N-linked site
classification to exact glycan identification over 2,357 candidates; the
benchmark is model-agnostic and we provide reference implementations
rather than claiming state-of-the-art models.  Third, we perform a
\textbf{systematic investigation of cellular context} using single-cell
enzyme expression data, providing the first quantitative evidence that
current enzyme--glycan annotations are insufficient to bridge the
determination gap. Together, these analyses delineate a boundary between
sequence-predictable structural classes and environment-dependent glycan
identity, and identify specific data gaps that must be filled to advance the
field.

% ══════════════════════════════════════════════════════════════════
\section{Materials and methods}

\subsection{Knowledge graph construction}

The glycoMusubi ETL pipeline comprises four automated stages: \textbf{download},
\textbf{clean}, \textbf{build} and \textbf{validate} (Supplementary Fig.~S1).
Data are fetched from six public databases through rate-limited API clients
with configurable retry logic and exponential backoff.

\textbf{Data sources.} GlyGen (v1.28, accessed January 2025) provides
glycosyltransferase annotations and glycan master lists. GlyTouCan (v3.2.1,
accessed January 2025) provides glycan structures in WURCS notation via the
GlyCosmos SPARQList API. UniProt (release 2024\_06) provides protein annotations
including disease associations, genetic variants and glycosylation site features.
ChEMBL (v34, accessed February 2025) provides glycosyltransferase inhibitor
activity data. PTMCode (v3) provides post-translational modification site-level
crosstalk network data. GlyConnect (v2.0, accessed March 2025) provides
site-specific glycan assignments with protein and position annotation. Reactome
(release 88, accessed February 2025) provides curated biosynthetic reaction
pathways linking enzymes to glycan products.

\textbf{Schema enforcement.} The KG schema defines 10 node types and 14
relation types (Supplementary Table~S1). Node identifiers are validated against
source-specific patterns (GlyTouCan: \texttt{\^{}G\textbackslash d\{5\}[A-Z]\{2\}\$};
UniProt: Swiss-Prot accession format; ChEMBL: \texttt{\^{}CHEMBL\textbackslash d+\$}).
Edges are constrained to declared source--target type pairs. Validation checks
include orphan node detection, duplicate edge removal, self-loop elimination
and cycle detection in acyclic relations.

\textbf{Glycosylation site processing.} UniProt glycosylation sites are
filtered to retain entries with experimental or strong sequence-similarity
evidence (evidence codes ECO:0000269 [experimental], ECO:0000303 [traceable
author statement] and ECO:0000250 [sequence similarity to experimentally
characterised orthologues]), discarding computational predictions and
low-confidence annotations. Site positions are normalised to canonical isoform
coordinates using the standardised format
\texttt{SITE::<UniProtID>::<Position>::<Residue>}.

\textbf{PTM crosstalk integration.} PTMCode crosstalk edges (minimum confidence
score 0.30) link pairs of modification sites with evidence of functional
interaction across six PTM types (phosphorylation, methylation, acetylation,
ubiquitination, SUMOylation and glycosylation).

\subsection{Knowledge graph embedding model}

To obtain glycan representations for structural clustering and retrieval, we
trained a knowledge graph embedding model (GlycoKGNet) on the glycoMusubi graph.

\textbf{Modality-specific encoders.} Glycan nodes are encoded by a
GlycanTreeEncoder that operates on the WURCS-parsed branching tree structure
through bottom-up tree message-passing layers (3 layers, hidden dimension 256)
with child attention and GRU state updates, followed by top-down refinement and
attention pooling, producing 256-dimensional glycan embeddings. The tree encoder
explicitly represents monosaccharide identity, linkage type (position and
anomeric configuration), and branching topology---features that distinguish it
from sequence-based approaches such as SweetNet~\cite{burkholz2022}
(Supplementary Section~S12 provides a detailed architecture comparison and
ablation study). The embedding dimension (256) was selected via grid search over
$\{64, 128, 256, 512\}$ on validation MRR (Supplementary Table~S16). Protein
nodes are encoded by projecting pre-computed ESM-2~\cite{lin2023} per-residue
embeddings (650M parameter model, frozen) through a site-aware MLP that
incorporates glycosylation site positional information. Disease and pathway nodes
are encoded by projecting PubMedBERT text embeddings (110M parameters, frozen)
through learned MLPs.

\textbf{Link prediction training.} The model is trained on the KG link prediction
task using type-constrained negative sampling (64 negatives per positive) with
self-adversarial weighting. The scoring function combines DistMult and RotatE
sub-scorers with per-relation adaptive weights. The training loss is binary
cross-entropy with structural contrastive regularisation (InfoNCE,
$\tau = 0.07$) and L2 regularisation. Optimisation uses Adam with cosine
annealing and mixed-precision training.

\textbf{Inductive evaluation.} We train in the inductive setting, holding out
15\% of entities (including their edges) during training and evaluating the
model's ability to embed and score triples involving unseen entities. The
inductive filtered MRR on the held-out set is 0.207 (Hits@10 = 0.522),
demonstrating that the modality-specific encoders generalise to entities not
seen during training. While the overall inductive MRR is lower than transductive
baselines (e.g.\ TransE MRR = 0.474; Supplementary Table~S4), GlycoKGNet
substantially outperforms TransE on the biologically critical
\texttt{has\_glycan} relation (3.5$\times$ improvement; Supplementary Table~S5).
This is essential for the glycan prediction application, where the model must
embed glycans that may not be present in the training KG.

\textbf{Preventing information leakage.} Glycan hierarchy edges (parent\_of,
child\_of, subsumes, subsumed\_by) could leak structural information across
splits. To prevent this, we remove all hierarchy edges involving held-out
entities from the training graph. The protein-stratified split for downstream
tasks further ensures that no protein or its associated sites appear in both
training and test sets. We verified that removing all glycan hierarchy edges (parent\_of, child\_of,
subsumes, subsumed\_by) from both KG embedding training and the subsequent
KMeans clustering input changes Task~5 H@5 by less than 0.02, confirming
that the downstream results are not driven by hierarchy-mediated leakage
(Supplementary Table~S13).

\subsection{Hierarchical prediction tasks}

\textbf{N-linked site classification.} A two-layer MLP (256 $\to$ 128 $\to$ 1)
takes ESM-2 per-residue embeddings at asparagine positions as input and predicts
N-linked glycosylation probability. Training uses binary cross-entropy with
class-balanced sampling. Evaluation on held-out proteins reports F1, precision,
recall and AUC-ROC.

\textbf{Site ranking.} Given a protein with multiple confirmed glycosylation
sites, a pairwise ranking model scores sites by glycosylation propensity using
ESM-2 embeddings with a margin-based ranking loss. Evaluation reports Recall@$k$
and MRR among known sites per protein.

\textbf{Glycosylation type prediction.} A fraction regression model predicts the
proportional composition of high-mannose, complex, hybrid and unoccupied states
at each site, using ESM-2 embeddings. The dominant type is taken as the argmax
prediction; accuracy is computed across four classes.

\textbf{Structural family classification.} An 8-class MLP classifier assigns
each glycan observation to a structural family derived from glycan taxonomy.
Input features are ESM-2 per-residue embeddings; evaluation reports top-1 and
top-2 accuracy.

\textbf{Structural cluster prediction.} Glycan embeddings from the trained KG
model are clustered by KMeans ($K = 20$, selected by silhouette analysis over
$K \in \{5, 10, 20, 50, 100\}$; Supplementary Fig.~S2) and the retrieval model's
top-$k$ cluster assignments (by majority vote among the $k$ nearest glycan
embeddings) are evaluated against the true glycan's cluster. The retrieval model
uses cosine similarity between site-level protein representations (ESM-2 MLP) and
glycan embeddings in the shared 256-dimensional space. The determination boundary
is robust to cluster granularity: the gap between cluster-level H@5 and exact-ID
MRR exceeds 5-fold for all tested $K$ values (Supplementary Table~S12).

\textbf{Exact glycan identification.} The same retrieval model ranks all 2,357
N-linked glycan candidates by cosine similarity to the query protein/site
representation. Two variants are evaluated:
site-level (using per-residue ESM-2 features at the glycosylation position,
ranking against 2,357 distinct N-linked glycan candidates) and protein-level
(using mean-pooled ESM-2 features, ranking against all 3,345 glycans including
O-linked structures to assess cross-type generalisation). Evaluation follows the filtered ranking
protocol~\cite{bordes2013}, reporting MRR, Hits@1, Hits@3, Hits@10 and mean
rank.

\subsection{Cell-type conditioned glycan prediction}

\textbf{Data preparation.} Cell-type enzyme expression profiles are obtained
from the Tabula Sapiens atlas~\cite{tabulasapiens2022}: for each of 748 cell
types and 339 glycosylation-related enzyme genes (defined as all genes annotated
with glycosyltransferase or glycosidase activity in GlyGen, filtered to those
present in the Tabula Sapiens expression matrix), the mean log-normalised
expression value is recorded. Enzyme-to-glycan biosynthetic edges from glycoMusubi
map each enzyme gene to the set of glycan structures it is known to produce.

\textbf{Within-cluster enzyme scoring.} For each cell type $c$ and glycan $g$ in
cluster $k$, the enzyme score is:
\begin{equation}
s(g, c) = \sum_{e \in \text{enzymes}(g)} \text{expr}(e, c)
\end{equation}
where $\text{enzymes}(g)$ is the set of enzyme genes linked to glycan $g$ in the
KG and $\text{expr}(e, c)$ is the expression level of enzyme $e$ in cell type
$c$. Glycans with no enzyme links receive a score of zero.

\textbf{Evaluation modes.} Four modes are evaluated within each cluster:
(1)~\textbf{Marginal}: enzyme scores averaged over all 748 cell types;
(2)~\textbf{Oracle}: the cell type giving the best rank for the true glycan
(upper bound on cell-type selection);
(3)~\textbf{Popularity}: ranking by training-set glycan frequency;
(4)~\textbf{Enzyme + popularity}: linear blend
$s_{\text{blend}} = \alpha \cdot s_{\text{enzyme}} + (1 - \alpha) \cdot s_{\text{pop}}$
with $\alpha$ selected on a validation set.

\textbf{Learned scorer.} A CellTypeGlycanScorer neural network maps enzyme
expression vectors (339 dimensions) through a two-layer MLP with LayerNorm and
GELU activation, and glycan embeddings (256 dimensions from the KG model)
through a parallel projection, producing normalised vectors whose inner product
(with learnable temperature and per-glycan bias) scores each glycan. Training
uses InfoNCE loss within each cluster, with cell-type expression averaged across
the training set (since the true cell type per sample is unknown).

\textbf{Cascade estimation.} The full-ranking cascade MRR can be approximated as
the product of cluster-level Hits@1 (0.166) and within-cluster MRR under an
independence assumption. However, since cluster prediction errors systematically
bias the within-cluster candidate set, we additionally evaluate the cascade
end-to-end: for each test sample, the model first predicts the cluster, then
ranks glycans within the predicted cluster. The end-to-end cascade MRR is 0.017,
compared to the independence-based estimate of 0.017, indicating that the two
stages are approximately independent in practice.\footnote{This approximate
independence holds across all five protein-stratified random splits: end-to-end
MRR ranges from 0.015 to 0.019, while independence estimates range from 0.015
to 0.018 (relative difference $<$15\% in all cases, well within the bootstrap
95\% CI width of $\pm$0.004 for cascade MRR).} This end-to-end result confirms
that neither stage alone is the bottleneck---both cluster prediction (H@1 = 0.166)
and within-cluster ranking (MRR = 0.105) require improvement.

\subsection{Statistical procedures}

All experiments use a single train/validation/test split stratified by protein
(no protein appears in multiple splits). To assess stability, we report 95\%
non-parametric percentile bootstrap confidence intervals (1,000 resamples of the
test set, sampling with replacement at the protein level to respect the
stratification) for all primary metrics (Table~\ref{tab:hierarchical}). Additionally, we verified that the key finding---the 10-fold gap between
cluster-level and exact-ID performance---is robust across five independent
protein-stratified random splits (Supplementary Table~S11).
Retrieval metrics are computed under the
filtered ranking protocol~\cite{bordes2013}, where known true associations are
excluded from the candidate ranking to avoid trivially inflated scores.
The two-layer boundary is validated by the consistent
performance gap across model architectures and feature representations: the
cluster-level advantage over exact-ID prediction persists for both site-level
(0.786 versus 0.066) and protein-level (0.786 versus 0.118) models, ruling out
model-specific artefacts. Because all comparisons are descriptive (performance
gaps across task granularities rather than competing models on the same task), we
do not apply multiple-comparison corrections; the consistency across architectures
serves as the primary evidence of robustness.

\textbf{Homology considerations.} The protein-stratified split prevents
the same protein from appearing in both training and test sets but does
not explicitly control for sequence homology between proteins across
splits. To quantify potential homology leakage, we computed pairwise
sequence identities between test and training proteins using MMseqs2
(sensitivity 7.5) and found that 8.3\% of test proteins share
$\geq$30\% sequence identity with at least one training protein. To
assess the impact, we repeated the Task~1 evaluation excluding these
proteins: F1 decreased from 0.924 to 0.908, indicating that homology
contributes modestly to Task~1 performance but does not fundamentally
alter the result. For Tasks~5--6 (retrieval), the effect is smaller
(MRR change $<$0.003) because the retrieval model operates in KG
embedding space rather than directly on sequence features
(Supplementary Table~S18). A fully homology-clustered split (e.g.\ at
30\% identity using MMseqs2) is desirable but would reduce the already
limited test set substantially; we leave this for future work with
larger glycoproteomics datasets.

\subsection{Implementation}

glycoMusubi is implemented in Python 3.11 using PyTorch and PyTorch Geometric.
Pre-computed ESM-2 per-residue embeddings and PubMedBERT text embeddings are
cached as \texttt{.pt} files. Configuration is managed through hierarchical YAML
files with OmegaConf. The full pipeline---ETL, embedding, prediction,
evaluation---is executable through scripted entry points. A Docker image with all
dependencies is provided for reproducible deployment.

% ══════════════════════════════════════════════════════════════════
\section{Results}

\subsection{Knowledge graph construction}

We constructed glycoMusubi through an automated four-stage ETL pipeline (download,
clean, build, validate) that integrates data from six public databases
(Fig.~\ref{fig:kg_benchmark}a). The resulting knowledge graph contains 78,263
nodes spanning 10 biological entity types (glycans, proteins, enzymes, diseases,
variants, compounds, PTM sites, motifs, reactions and cellular locations) and
2,508,960 edges spanning 14 relation types (Table~\ref{tab:stats}).

Identifier normalisation resolves conflicts across databases: UniProt isoform
accessions are mapped to canonical forms, GlyTouCan accessions are validated
against source patterns, and glycosylation sites are encoded as structured
identifiers enabling unambiguous cross-referencing.

\begin{table}[ht]
\centering
\caption{glycoMusubi statistics.}
\label{tab:stats}
\small
\begin{tabular}{lrl}
\toprule
\textbf{Entity type} & \textbf{Count} & \textbf{Primary source} \\
\midrule
Glycan & 58,628 & GlyTouCan \\
Protein & 4,372 & UniProt \\
Enzyme & 1,459 & GlyGen \\
Site & 10,218 & UniProt, PTMCode \\
Disease & 1,936 & UniProt \\
Compound & 746 & ChEMBL \\
Variant & 522 & UniProt \\
Motif & 181 & GlyTouCan \\
Reaction & 164 & Reactome \\
Cellular location & 37 & UniProt \\
\midrule
\textbf{Total nodes} & \textbf{78,263} & \\
\textbf{Total edges} & \textbf{2,508,960} & 14 relation types \\
\bottomrule
\end{tabular}
\end{table}

\subsection{Hierarchical glycan prediction benchmark}

We decomposed glycan prediction into a hierarchy of six tasks, each narrowing
the candidate space and increasing specificity
(Fig.~\ref{fig:kg_benchmark}b, Table~\ref{tab:hierarchical}).

\textbf{Task~1: N-linked site classification.} Using ESM-2~\cite{lin2023}
per-residue embeddings as input to a two-layer MLP, we achieve F1 = 0.924
(AUC-ROC = 0.825) on held-out proteins, confirming that the local sequence
context is a strong predictor of N-linked glycosylation potential.

\textbf{Task~2: Site ranking.} An ESM-2-based ranking model achieves
Recall@3 = 0.683 and MRR = 0.688, indicating that protein sequence features
capture site-level glycosylation propensity.

\textbf{Task~3: Glycosylation type prediction.} A fraction regression model
achieves accuracy = 0.844 across four categories (high-mannose, complex, hybrid,
unoccupied), demonstrating that the structural category of glycosylation is
partially encoded in the protein sequence.

\textbf{Task~4: Structural family (8 classes).} An ESM-2 MLP achieves top-2
accuracy = 0.858 and top-1 accuracy = 0.638, suggesting that protein sequence
reliably narrows the glycan to a small set of structural families.

\textbf{Task~5: Structural cluster ($K = 20$).} Clusters are defined by KMeans
over KG-derived glycan embeddings (mean cluster size $\sim$118 glycans). A
site-level retrieval model achieves Hits@5 = 0.786 and Hits@10 = 0.978,
indicating that protein sequence and site context predict the correct structural
neighbourhood with high accuracy. For comparison, a random baseline achieves H@5 = 0.25 (expected value for $K = 20$) and a popularity baseline achieves H@5 = 0.412, confirming that the retrieval model substantially outperforms non-trivial baselines at the cluster level.

\textbf{Task~6: Exact glycan identification.} Site-level retrieval ranks against
2,357 N-linked glycans, achieving MRR = 0.066 (Hits@10 = 0.145). Protein-level
retrieval ranks against 3,345 glycans, achieving MRR = 0.118 (Hits@10 = 0.228).
The approximately 10-fold gap between cluster-level performance (Task~5) and
exact identification (Task~6) reveals an empirical determination boundary.

\textbf{Sequence-based retrieval versus frequency baseline.} Notably, the
site-level retrieval model (MRR = 0.066) does not surpass the popularity
baseline (MRR = 0.100), which simply ranks glycans by their training-set
frequency. The protein-level variant (MRR = 0.118) modestly exceeds the
popularity baseline but remains far below cluster-level performance.
This indicates that, at exact-ID resolution, protein sequence alone does
not add sufficient discriminative information beyond glycan frequency
statistics. The popularity baseline's strength reflects the highly skewed
glycan frequency distribution (the 10 most common N-linked glycans account
for 34\% of all observations), which provides a strong prior that
sequence features cannot yet overcome at fine granularity.  This
observation reinforces our central conclusion: structural-class prediction
is tractable from sequence (Tasks~1--5), but exact glycan identity remains
an open problem requiring information sources beyond protein sequence
and currently available annotations.

\textit{Note on metrics.} Each task uses the metric most appropriate to its
formulation: classification tasks report F1 or accuracy, ranking tasks report
MRR, and retrieval tasks report Hits@$k$. To enable cross-task comparison, we
additionally report a normalised ``reduction in candidate space'' in
Supplementary Table~S6, confirming that the sharp performance drop occurs
specifically at the cluster-to-exact-ID transition regardless of metric choice.

\begin{table}[ht]
\centering
\caption{Hierarchical glycan prediction results. 95\% bootstrap confidence
intervals (1,000 protein-level resamples) are reported for primary metrics;
secondary metrics and baselines are point estimates.
\textsuperscript{\dag}Protein-level retrieval (3,345 candidates including O-linked glycans) achieves higher MRR than site-level retrieval (2,357 N-linked only). Two factors contribute: (i)~the candidate sets differ, making direct MRR comparison problematic; (ii)~mean-pooled ESM-2 features may act as implicit regularisation. We present both variants for completeness but caution against interpreting the protein-level result as evidence that site information is unhelpful; an ablation removing site-position information from the site-level model reduces MRR from 0.066 to 0.058 (Supplementary Table~S17), confirming that positional information provides a modest but real benefit within the same candidate set.}
\label{tab:hierarchical}
\small
\begin{tabular}{llllrlr}
\toprule
\textbf{Task} & \textbf{Candidates} & \textbf{Model} & \textbf{Metric} & \textbf{Value} & \textbf{Secondary} & \textbf{Value} \\
\midrule
N-linked classification & Binary & ESM-2 MLP & F1 & 0.924{\scriptsize$\pm$0.008} & AUC-ROC & 0.825 \\
Site ranking & Per protein & ESM-2 ranking & Recall@3 & 0.683{\scriptsize$\pm$0.015} & MRR & 0.688 \\
Glycosylation type & 4 classes & Frac. regression & Accuracy & 0.844{\scriptsize$\pm$0.012} & & \\
Structural family & 8 classes & ESM-2 MLP & Top-2 Acc & 0.858{\scriptsize$\pm$0.014} & Top-1 Acc & 0.638 \\
Structural cluster & 20 clusters & Site-level retr. & H@5 & 0.786{\scriptsize$\pm$0.018} & H@10 & 0.978 \\
Exact glycan (site) & 2,357 glycans & Site-level retr. & MRR & 0.066{\scriptsize$\pm$0.004} & H@10 & 0.145 \\
Exact glycan (prot.)\textsuperscript{\dag} & 3,345 glycans & Protein-level retr. & MRR & 0.118{\scriptsize$\pm$0.006} & H@10 & 0.228 \\
\midrule
\multicolumn{7}{l}{\textit{Baselines for exact glycan identification}} \\
Random baseline & 2,357 glycans & --- & MRR & 0.001 & H@10 & 0.004 \\
Popularity baseline & 2,357 glycans & Train frequency & MRR & 0.100 & H@10 & 0.193 \\
Seq.\ similarity (BLAST) & 2,357 glycans & Nearest neighbour & MRR & 0.043 & H@10 & 0.109 \\
\bottomrule
\end{tabular}
\end{table}

\textbf{O-linked glycan prediction as a control experiment.} The O-linked
prediction problem provides an instructive contrast. With a substantially smaller
candidate space (fewer distinct O-linked glycan structures; 112 candidates versus
2,357 for N-linked, evaluated on 9,288 test samples), site-level retrieval
achieves MRR~=~0.847 (H@1~=~0.775, H@3~=~0.918; Supplementary Table~S3),
compared to MRR~=~0.066 for N-linked glycans under the same model architecture.
This 13-fold performance difference arises despite using identical training
procedures and model capacity, isolating candidate-space size and glycan
diversity as the primary factors driving N-linked prediction difficulty rather
than inherent model limitations or unpredictability of glycosylation. Notably,
the O-linked result demonstrates that the retrieval architecture is capable of
near-perfect glycan identification when the candidate space is tractable.

\subsection{Cell-type enzyme expression analysis}

To investigate whether cellular context can bridge the determination gap, we
developed a cascade model combining sequence-based cluster prediction with
cell-type-specific enzyme expression scoring (Fig.~\ref{fig:twolayer}a). The
cascade decomposes glycan prediction as:
\begin{equation}
P(\text{glycan} \mid \text{protein, site, cell\_type}) = P(\text{cluster} \mid \text{protein, site}) \times P(\text{glycan} \mid \text{cluster, cell\_type})
\end{equation}

Using enzyme expression profiles from the Tabula Sapiens single-cell
atlas~\cite{tabulasapiens2022} (748 cell types, 339 glycosylation-related
enzyme genes) and enzyme-to-glycan biosynthetic edges from glycoMusubi, we scored
glycans within each predicted cluster by the expression levels of their
producing enzymes.

The results reveal a critical data gap (Table~\ref{tab:celltype}). Within the
cluster context (average $\sim$118 candidates), the popularity baseline achieves
MRR = 0.105, while the enzyme expression marginal achieves only MRR = 0.035 and
the oracle achieves MRR = 0.026---counter-intuitively lower than the marginal, a phenomenon we trace to the extreme sparsity of enzyme--glycan annotations (see Table~\ref{tab:celltype} footnote). The enzyme+popularity blend assigns optimal
weight $\alpha = 0.0$ to the enzyme component. A learned neural scorer achieves
MRR = 0.104---matching the popularity baseline.

The root cause is annotation sparsity: only 929 of 2,357 N-linked glycans
(39.4\%) have any enzyme link in glycoMusubi, and these are concentrated in a
single mega-cluster containing 1,843 glycans (78\% of all N-linked glycans).
Of the 20 structural clusters, only one has non-trivial enzyme coverage.
Crucially, the sparsity is non-uniform across enzyme categories: core pathway
enzymes (e.g.\ MGAT1, MGAT2, MAN1A1) cover 87\% of annotated glycans but are
ubiquitously expressed across cell types, providing no discriminative signal.
Terminal processing enzymes (e.g.\ ST3GAL1--6, FUT1--11, B4GALT1--4) are highly
cell-type-specific~\cite{chrysinas2024} but have enzyme--glycan annotations for
only 12\% of N-linked glycans in glycoMusubi. This precisely explains why cell-type
expression adds no predictive value: the informative enzymes lack glycan
annotations, while the well-annotated enzymes lack cell-type specificity
(Supplementary Table~S14).

\textbf{Glycosylation type and prediction difficulty.} Performance varies
substantially across structural categories. High-mannose glycans, which undergo
minimal Golgi processing, are predicted with MRR = 0.142 (site-level), while
complex-type glycans---subject to extensive and variable terminal
modifications---achieve only MRR = 0.038. Hybrid-type glycans fall between
(MRR = 0.071). Paucimannose glycans, which undergo trimming but minimal
elongation, are also well-predicted (MRR = 0.194, 42 candidates). This
gradient confirms that prediction difficulty correlates with the degree of
environment-dependent enzymatic processing (Supplementary Table~S15).

\begin{table}[ht]
\centering
\caption{Cell-type conditioned glycan prediction.
\textsuperscript{$\ddagger$}The oracle selects, for each test sample, the cell type maximising the rank of the true glycan among \emph{enzyme-annotated} glycans only.  Because 60.6\% of test glycans lack any enzyme annotation and therefore receive a score of zero regardless of cell type, the oracle cannot improve their rank.  The marginal averages over all cell types, which smooths enzyme scores and incidentally provides non-zero scores to a broader set of glycans through correlated expression patterns. Consequently, the oracle's restriction to per-sample optimisation over a sparse signal yields lower aggregate MRR than the marginal's averaging effect.  This counter-intuitive result is itself diagnostic of the annotation sparsity problem.}
\label{tab:celltype}
\small
\begin{tabular}{llrrrp{3.5cm}}
\toprule
\textbf{Method} & \textbf{Scope} & \textbf{MRR} & \textbf{H@10} & \textbf{H@50} & \textbf{Note} \\
\midrule
Popularity baseline & Full (2,357) & 0.100 & 0.193 & 0.439 & Train frequency \\
Enzyme marginal & Full (2,357) & 0.001 & 0.000 & 0.000 & All cell types averaged \\
\midrule
Popularity & Within-cluster & 0.105 & 0.207 & 0.469 & Avg $\sim$118 candidates \\
Enzyme marginal & Within-cluster & 0.035 & 0.034 & 0.034 & Avg $\sim$118 candidates \\
Enzyme oracle\textsuperscript{$\ddagger$} & Within-cluster & 0.026 & 0.026 & 0.029 & Best cell type per sample \\
Enzyme + pop. & Within-cluster & 0.105 & 0.207 & 0.469 & Optimal $\alpha = 0.0$ \\
Learned scorer & Within-cluster & 0.104 & 0.205 & 0.467 & InfoNCE-trained \\
\midrule
Site-level retrieval & Full (2,357) & 0.066 & 0.145 & --- & Sequence only (ref.) \\
Protein-level retr. & Full (3,345) & 0.118 & 0.228 & --- & Sequence only (ref.) \\
\bottomrule
\end{tabular}
\end{table}

\subsection{Systematic negative results}

Several additional approaches failed to improve prediction beyond baselines
(Table~\ref{tab:negative}). Monosaccharide composition prediction from protein
sequence yielded $R^2 < 0$ for all sugar types, confirming no learnable
relationship. Biosynthetic pathway databases (glycowork subgraphs,
WURCS-based reactions) did not improve over the popularity baseline. RNA
expression levels from the Human Protein Atlas~\cite{uhlen2015} achieved
MRR = 0.001, far below the popularity baseline. Enzyme pathway scoring as an
auxiliary retrieval signal yielded optimal blending weight $\alpha = 1.0$ for the
retrieval score alone.

\begin{table}[ht]
\centering
\caption{Negative results summary.}
\label{tab:negative}
\small
\begin{tabular}{llcrp{3.5cm}}
\toprule
\textbf{Approach} & \textbf{Data source} & \textbf{Metric} & \textbf{Value} & \textbf{Interpretation} \\
\midrule
Composition pred. & ESM-2 $\to$ monosaccharides & $R^2$ & $< 0$ & Not learnable from sequence \\
Reaction network & glycowork & MRR & $\leq$ pop. & Annotations too sparse \\
Tissue expression & HPA RNA consensus & MRR & 0.001 & Bulk tissue insufficient \\
Enzyme pathway & GlyGen enzyme edges & $\alpha^*$ & 1.0 & No added information \\
Cell-type enzyme & Tabula Sapiens & MRR & 0.035 & Annotations too sparse \\
\bottomrule
\end{tabular}
\end{table}

\subsection{Two-layer determination model}

The hierarchical benchmark and negative-result analyses together support an
empirical determination boundary under currently available data, modelled as two distinct layers
(Fig.~\ref{fig:twolayer}b).

\textbf{Layer~1: Sequence-determined.} Protein sequence and site context
strongly constrain the structural class of glycosylation (Tasks~1--5;
F1 = 0.924 for site classification, H@5 = 0.786 for structural cluster). The
biological basis is glycosyltransferase substrate
specificity~\cite{schjoldager2020,narimatsu2019}.

\textbf{Layer~2: Environment-determined.} Within a structural class, exact
glycan identity depends on cellular factors: enzyme expression, nucleotide sugar
availability, Golgi transit time and competing biosynthetic
pathways~\cite{stanley2022,varki2022essentials}. This corresponds to Task~6,
where performance drops sharply (MRR = 0.066--0.118). Our cell-type analysis
shows that current annotations are too sparse to bridge this gap computationally
(39\% glycan coverage, concentrated in one cluster).

The boundary is not a limitation of the models used: the 10-fold gap persists
across model architectures (site-level versus protein-level retrieval) and
feature representations (ESM-2 MLP versus embedding-based retrieval), suggesting
it reflects properties of glycosylation biology that are robust across the models and representations tested here, though we cannot exclude that future methods with richer input data may narrow this gap
(Fig.~\ref{fig:twolayer}c).

% ══════════════════════════════════════════════════════════════════
\section{Discussion}

Our hierarchical benchmark provides the first systematic quantification of what
protein sequence can and cannot predict about glycosylation. The high accuracy of
structural-class prediction (H@5 = 0.786 for $K = 20$ clusters) is consistent
with the known specificity of glycosyltransferases for particular acceptor
sequences~\cite{schjoldager2020} and suggests that sequence-based tools can
reliably predict \textit{what type} of glycan modification occurs at a given
site. This level of prediction is sufficient for many applications in structural
biology, molecular dynamics simulation and rational protein engineering.

We emphasise that the primary contribution of this work is the
integrated resource and benchmark rather than the specific models used.
GlycoKGNet and the retrieval models serve as reference implementations
to establish baseline performance on the benchmark tasks; we expect
that future methods---particularly those incorporating structural
predictions from AlphaFold or spatial glycomics data---will improve upon
these baselines, and the hierarchical benchmark provides a standardised
framework for measuring such progress.

The sharp drop to MRR = 0.066--0.118 for exact glycan identification means that
predicting the precise glycan structure requires information beyond the protein
sequence. Our systematic investigation of candidate information
sources---single-cell enzyme expression, tissue expression profiles, reaction
networks and composition---found that none currently bridges this gap.

The failure of enzyme expression is particularly informative. The Tabula Sapiens
atlas provides the most detailed available view of cell-type-specific enzyme
expression, yet it adds no predictive signal. This is not because the biological
hypothesis is wrong---glycan biosynthesis is indeed determined by enzyme
expression---but because the enzyme--glycan annotation mapping is critically
incomplete. Only 929 of 2,357 N-linked glycans (39.4\%) have any enzyme link in
current databases, concentrated in a single structural cluster. This annotation
gap, rather than model capacity, is the primary bottleneck.

The O-linked prediction result (MRR = 0.847) provides important context. The
high accuracy for O-linked glycans---which have a smaller candidate
space---suggests that the N-linked challenge is largely one of candidate-space
size and glycan diversity rather than inherent unpredictability.

\textbf{Comparison with prior work.} Our work is most directly comparable to
InSaNNE~\cite{kellman2024}, which predicts glycan structures from protein
sequence features using a recurrent neural network. InSaNNE achieves AUC = 0.92
on binary presence/absence classification across 199 glycans and 57 glycosites.
Our approach differs in three ways: (i)~we evaluate retrieval ranking (MRR) over
2,357 glycans rather than binary classification over 199, reflecting the full
scale of the problem; (ii)~we decompose the problem into a hierarchy of six
tasks, revealing where prediction succeeds and fails; and (iii)~we systematically
investigate cellular context as an additional information source. The
determination boundary we identify---reliable structural-class prediction but
poor exact-ID prediction---is consistent with Glycoimpact's `bounded
biosynthesis' framework~\cite{glycoimpact2024}, which showed that protein
sequence constrains but does not determine the glycan.

Yadav \textit{et al.}~\cite{yadav2022} modelled Golgi glycan processing as an
optimal information-coding channel, predicting that structural diversity is
bounded by enzyme kinetics and compartmental organisation. Our empirical
determination boundary provides quantitative support for this theoretical
framework: the 10-fold gap between cluster-level and exact-ID prediction is broadly consistent with the information-theoretic constraints imposed by enzymatic diversity in the Golgi, though the precise magnitude of this gap may partly reflect current data limitations.

Chrysinas \textit{et al.}~\cite{chrysinas2024} used the same Tabula Sapiens
atlas to analyse glycosylation enzyme expression, finding that terminal
processing enzymes (e.g.\ sialyltransferases, fucosyltransferases) are highly
cell-type-specific while core pathway enzymes (e.g.\ MGAT1, MGAT2) are
ubiquitously expressed. Their qualitative observation aligns with our
quantitative finding that enzyme expression adds no predictive signal: the
cell-type-variable enzymes are precisely those with the sparsest glycan
annotations (39\% coverage), preventing their discriminative use.

SweetNet~\cite{burkholz2022}, GIFFLAR~\cite{thomes2024} and
GlycanML~\cite{glyml2024} predict glycan properties from structures (the forward
direction) but do not address the inverse problem. CoNglyPred~\cite{congly2025}
predicts N-linked glycosylation sites using ESM-2 and AlphaFold2 features,
complementing our Task~1 (F1 = 0.924) but not addressing glycan identity. The
glycoMusubi resource integrates data from six databases that prior approaches used
individually, enabling multi-hop queries that no single-database approach
supports.

\textbf{Limitations.} The benchmark is based on GlyConnect's (v2.0, accessed
March 2025) site-specific glycan assignments, which are biased toward
well-studied proteins (Gini coefficient of the protein distribution = 0.72;
the top 5\% of proteins contribute 38\% of test samples). ESM-2 per-residue embeddings are available for only 1,611
of 9,138 proteins (17.6\%) in the KG, restricted to those with confirmed
glycosylation sites in GlyConnect; all benchmark proteins have embeddings, but
the KG embedding model's protein coverage is limited. The structural clustering
depends on KG-derived glycan embeddings. The cell-type analysis assumes that
enzyme expression at the RNA level reflects enzymatic activity---a simplification
given that post-translational regulation of glycosyltransferases is known to
occur~\cite{schjoldager2020}. The current KG does not include mass spectrometry
glycan identification data.

\textbf{Future directions.} Closing the determination gap requires parallel
efforts in data and methods. Systematic characterisation of enzyme--glycan
biosynthetic relationships---particularly for complex-type N-linked
glycans---would directly increase coverage. Integrating spatial glycomics and
mass-spectrometry-based glycan identification could substantially improve
prediction at the exact-glycan level. The hierarchical benchmark provides a
standardised evaluation framework for measuring progress.

% ══════════════════════════════════════════════════════════════════
\section*{Data availability}

glycoMusubi integrates publicly available data from GlyGen
(\url{https://data.glygen.org}), GlyTouCan (\url{https://glytoucan.org}),
UniProt (\url{https://www.uniprot.org}), ChEMBL
(\url{https://www.ebi.ac.uk/chembl}), PTMCode (\url{https://ptmcode.embl.de})
and GlyConnect (\url{https://glyconnect.expasy.org}). The constructed knowledge graph, trained embeddings and benchmark results are
available as a reviewer-accessible pre-release on Zenodo
(DOI:~10.5281/zenodo.15443928; embargo lifted upon acceptance).

\section*{Funding}

This work was supported by JSPS KAKENHI (Grant Number JP24K02954) and the
Institute for Glyco-core Research (iGCORE), Nagoya University.

% ══════════════════════════════════════════════════════════════════
% Figures
% ══════════════════════════════════════════════════════════════════

% Figure 1: glycoMusubi resource and hierarchical benchmark overview
% Publication-quality TikZ — Bioinformatics / Nature style
\begin{figure*}[!ht]
\centering

% ═══════════════════════════════════════════════════════════════════
% Semantic colour palette (restrained: 3 hues + neutrals)
% ═══════════════════════════════════════════════════════════════════
\definecolor{cPrimary}{HTML}{2b6a99}     % muted steel blue — infrastructure
\definecolor{cPrimaryLight}{HTML}{d6e6f0} % very pale blue tint
\definecolor{cAccent}{HTML}{c04e3f}       % muted warm red — boundary / alert
\definecolor{cAccentLight}{HTML}{f5dcd8}  % pale red tint
\definecolor{cDark}{HTML}{1d1d2c}         % near-black for text & outlines
\definecolor{cMid}{HTML}{5a6068}          % medium gray for secondary text
\definecolor{cLight}{HTML}{b0b5bb}        % light gray for hairlines
\definecolor{cFaint}{HTML}{e8eaed}        % very faint gray for grid
\definecolor{cWhite}{HTML}{ffffff}
% KG node tones — two-tone scheme: primary entities vs. secondary
\definecolor{cNodeA}{HTML}{2b6a99}        % primary entity (protein, enzyme, glycan)
\definecolor{cNodeB}{HTML}{6b97b5}        % secondary entity

% ═══════════════════════════════════════════════════════════════════
% Global style definitions
% ═══════════════════════════════════════════════════════════════════
\begin{tikzpicture}[
    x=1cm, y=1cm,
    % ── reusable styles ──
    panellabel/.style={font=\fontsize{10pt}{12pt}\selectfont\bfseries,
                       text=cDark, anchor=north west},
    seclabel/.style={font=\fontsize{6.5pt}{8pt}\selectfont,
                     text=cMid, anchor=south},
    srcbox/.style={draw=cLight, rounded corners=2pt, fill=cPrimaryLight,
                   minimum width=2.0cm, minimum height=0.48cm,
                   font=\fontsize{6pt}{7pt}\selectfont\bfseries,
                   text=cPrimary, line width=0.35pt, inner sep=0pt},
    etlbox/.style={draw=cPrimary!70!black, rounded corners=2pt,
                   fill=cPrimary, text=white,
                   minimum width=1.6cm, minimum height=0.48cm,
                   font=\fontsize{6pt}{7pt}\selectfont\bfseries,
                   line width=0.35pt, inner sep=0pt},
    kgnode/.style={circle, minimum size=0.78cm,
                   font=\fontsize{4pt}{5pt}\selectfont\bfseries,
                   text=white, inner sep=0pt,
                   line width=0pt},
    thinedge/.style={cLight, line width=0.3pt},
    connector/.style={-{Stealth[length=1.4mm, width=0.9mm]},
                      cLight, line width=0.4pt},
    thickconn/.style={-{Stealth[length=1.6mm, width=1mm]},
                      cMid, line width=0.55pt},
    % bar chart
    barL1/.style={fill=cPrimary},
    barL2/.style={fill=cAccent},
    ticklabel/.style={font=\fontsize{5.5pt}{7pt}\selectfont, text=cMid},
    metriclabel/.style={font=\fontsize{5.5pt}{7pt}\selectfont, text=cDark},
    tasklabel/.style={font=\fontsize{6pt}{7pt}\selectfont, text=cDark},
    subtasklabel/.style={font=\fontsize{5pt}{6pt}\selectfont, text=cLight},
]

% ─────────────────────────────────────────────────────────────────
% PANEL (a) — top half: resource construction
% ─────────────────────────────────────────────────────────────────
\node[panellabel] at (-0.2, 7.0) {\textsf{a}};

% ── Section labels ──
\node[seclabel] at (1.0, 6.65) {Data sources};
\node[seclabel] at (3.7, 6.65) {ETL pipeline};
\node[seclabel] at (8.5, 6.65) {Knowledge graph};

% ── Thin horizontal rule beneath section labels ──
\draw[cFaint, line width=0.3pt] (-0.1, 6.55) -- (11.3, 6.55);

% ── Source column (left) ──────────────────────────────────────
\foreach \i/\name in {0/GlyGen, 1/GlyTouCan, 2/UniProt,
                       3/ChEMBL, 4/PTMCode, 5/Reactome}{
    \node[srcbox] (src\i) at (1.0, {6.1 - \i*0.62}) {\name};
}

% ── Merge connector: sources → bus → ETL ──────────────────────
% Vertical collector line
\draw[cLight, line width=0.6pt] (2.55, 2.95) -- (2.55, 6.1);
% Horizontal stubs from each source
\foreach \i in {0,...,5}{
    \draw[cLight, line width=0.3pt] (src\i.east) -- (2.55, {6.1 - \i*0.62});
}
% Trunk arrow → ETL
\draw[thickconn] (2.55, 4.52) -- (2.75, 4.52);

% ── ETL pipeline (centre) ────────────────────────────────────
\foreach \j/\stage in {0/Download, 1/Clean, 2/Build, 3/Validate}{
    \node[etlbox] (etl\j) at (3.7, {5.85 - \j*0.70}) {\stage};
}
% Stage arrows
\foreach \j in {0,1,2}{
    \pgfmathtruncatemacro{\jnext}{\j+1}
    \draw[connector, cPrimary!60] (etl\j) -- (etl\jnext);
}
% Enclosing hairline box
\draw[cLight, rounded corners=3pt, line width=0.3pt, dashed]
    (2.72, 3.28) rectangle (4.68, 6.20);

% ── Output arrow: ETL → KG ───────────────────────────────────
\draw[thickconn] (4.68, 4.74) -- (5.55, 4.74);

% ── Knowledge graph (right, focal element) ────────────────────
% 10 nodes on an ellipse — centre (8.5, 4.2), rx=2.35, ry=2.05
\pgfmathsetmacro{\cx}{8.5}
\pgfmathsetmacro{\cy}{4.30}
\pgfmathsetmacro{\rx}{2.35}
\pgfmathsetmacro{\ry}{2.05}

% Pre-draw edges (behind nodes)
% Define node angles: glycan=90, protein=54, enzyme=18, site=-18,
%   disease=-54, compound=-90, motif=-126, reaction=-162,
%   variant=162, cell_loc=126
\foreach \angA/\angB in {90/54, 18/90, 54/-54, 54/-18, 54/162,
                          18/-90, 90/-126, 18/-162, 54/126}{
    \draw[cFaint, line width=0.25pt]
        ({\cx+\rx*cos(\angA)}, {\cy+\ry*sin(\angA)}) --
        ({\cx+\rx*cos(\angB)}, {\cy+\ry*sin(\angB)});
}

% Draw nodes
\foreach \ang/\col/\lab in {
    90/cNodeA/glycan,
    54/cNodeA/protein,
    18/cNodeA/enzyme,
    -18/cNodeB/site,
    -54/cNodeB/disease,
    -90/cNodeB/compound,
    -126/cNodeB/motif,
    -162/cNodeB/reaction,
    162/cNodeB/variant,
    126/cNodeB/cell loc.%
}{
    \node[kgnode, fill=\col] at ({\cx+\rx*cos(\ang)}, {\cy+\ry*sin(\ang)}) {\lab};
}

% ── Scale summary (elegant single line beneath graph) ─────────
\node[font=\fontsize{5.5pt}{7pt}\selectfont, text=cMid, anchor=north]
    at (\cx, {\cy - \ry - 0.40})
    {78,263 nodes\enspace\textbullet\enspace 2.5\,M edges\enspace\textbullet\enspace
     10 types\,/\,14 relations};

% ── Thin separator between panels ─────────────────────────────
\draw[cFaint, line width=0.25pt] (-0.1, 1.65) -- (11.3, 1.65);

% ─────────────────────────────────────────────────────────────────
% PANEL (b) — bottom half: hierarchical benchmark
% ─────────────────────────────────────────────────────────────────
\node[panellabel] at (-0.2, 1.40) {\textsf{b}};

% Layout: bars span x = [2.0, 9.0] → 7cm = performance 0–1
\pgfmathsetmacro{\bL}{2.0}   % bar left edge
\pgfmathsetmacro{\bW}{7.0}   % bar total width (=1.0)
\pgfmathsetmacro{\bH}{0.32}  % bar height
\pgfmathsetmacro{\bG}{0.60}  % bar gap (centre-to-centre spacing)

% ── Faint grid ──
\foreach \g in {0, 0.2, 0.4, 0.6, 0.8, 1.0}{
    \draw[cFaint, line width=0.2pt]
        ({\bL + \g*\bW}, -3.55) -- ({\bL + \g*\bW}, 1.15);
}

% ── Bars ──
% Macro: \drawbar{index}{task}{subtitle}{metric}{value}{style}
% y = 0.85 - index * bG
\newcommand{\drawbar}[6]{%
    \pgfmathsetmacro{\yc}{0.85 - #1*\bG}%
    \pgfmathsetmacro{\yt}{\yc + \bH/2}%
    \pgfmathsetmacro{\yb}{\yc - \bH/2}%
    \pgfmathsetmacro{\xr}{\bL + #5*\bW}%
    \fill[#6] (\bL, \yb) rectangle (\xr, \yt);%
    % task label (left)
    \node[tasklabel, anchor=east] at ({\bL - 0.12}, \yc) {#2};%
    % subtitle
    \node[subtasklabel, anchor=east] at ({\bL - 0.12}, {\yc - 0.17}) {#3};%
    % metric (right of bar)
    \node[metriclabel, anchor=west] at ({\xr + 0.10}, \yc) {#4};%
}

\drawbar{0}{N-linked classification}{Binary}{F1 = 0.924}{0.924}{barL1}
\drawbar{1}{Site ranking}{Per protein}{R@3 = 0.683}{0.683}{barL1}
\drawbar{2}{Structural family}{8 classes}{Top-2 = 0.858}{0.858}{barL1}
\drawbar{3}{Structural cluster}{20 clusters}{H@5 = 0.786}{0.786}{barL1}

% ── Determination boundary ──
\pgfmathsetmacro{\bndY}{0.85 - 3.5*\bG}
% Accent-tinted band
\fill[cAccentLight, opacity=0.35]
    (-0.1, {\bndY - 0.12}) rectangle (11.3, {\bndY + 0.12});
% Hairline
\draw[cAccent, line width=0.5pt] (\bL - 0.3, \bndY) -- ({\bL + \bW + 0.3}, \bndY);
% Label
\node[font=\fontsize{5pt}{6pt}\selectfont\itshape, text=cAccent,
      anchor=east]
    at ({\bL + \bW + 0.3}, {\bndY + 0.18}) {determination boundary};

% ── Below-boundary bar (accent colour) ──
\drawbar{4}{Exact glycan ID}{2{,}357 glycans}{MRR = 0.066}{0.066}{barL2}

% ── X-axis ──
\pgfmathsetmacro{\axY}{0.85 - 4*\bG - \bH/2 - 0.15}
\draw[cMid, line width=0.35pt] (\bL, \axY) -- ({\bL+\bW}, \axY);
\foreach \t/\l in {0/0, 0.2/0.2, 0.4/0.4, 0.6/0.6, 0.8/0.8, 1.0/1.0}{
    \draw[cMid, line width=0.25pt]
        ({\bL+\t*\bW}, \axY) -- ({\bL+\t*\bW}, {\axY - 0.06});
    \node[ticklabel, anchor=north] at ({\bL+\t*\bW}, {\axY - 0.08}) {\l};
}
\node[font=\fontsize{6pt}{7pt}\selectfont, text=cMid, anchor=north]
    at ({\bL + \bW/2}, {\axY - 0.38}) {Performance};

% ── Specificity arrow (left margin, spanning all bars) ──
\pgfmathsetmacro{\specTop}{0.85 + \bH/2}
\pgfmathsetmacro{\specBot}{0.85 - 4*\bG - \bH/2}
\draw[-{Stealth[length=1.2mm, width=0.8mm]}, cLight, line width=0.35pt]
    (0.35, \specTop) -- (0.35, \specBot);
\node[font=\fontsize{4.5pt}{5.5pt}\selectfont, text=cLight,
      rotate=90, anchor=south] at (0.20, {(\specTop+\specBot)/2})
    {increasing specificity};

\end{tikzpicture}

\caption{\textbf{glycoMusubi resource and hierarchical benchmark overview.}
\textbf{(a)}~Data integration pipeline: six public databases are processed
through an automated four-stage ETL pipeline to construct a heterogeneous
knowledge graph with 78{,}263 nodes (10~types) and 2{,}508{,}960 edges
(14~relation types).
\textbf{(b)}~Hierarchical prediction benchmark overview: performance across five
tasks of increasing specificity.  The determination boundary marks the sharp
performance drop from structural-class prediction (H@5\,=\,0.786) to exact
glycan identification (MRR\,=\,0.066). Fig.~\ref{fig:twolayer}c expands this
view with all seven tasks annotated by Layer~1/2 assignment.}
\label{fig:kg_benchmark}
\end{figure*}


% Figure 2: Two-layer determination model and cell-type analysis (TikZ)
\begin{figure}[!ht]
\centering

% ── Grayscale palette ──
\definecolor{cL1}{gray}{0.35}       % Layer 1: dark gray
\definecolor{cL1light}{gray}{0.92}  % Layer 1 background
\definecolor{cL1mid}{gray}{0.45}    % Layer 1 mid-tone
\definecolor{cL1dark}{gray}{0.20}   % Layer 1 darkest
\definecolor{cL2}{gray}{0.50}       % Layer 2: medium gray
\definecolor{cL2light}{gray}{0.92}  % Layer 2 background
\definecolor{cL2mid}{gray}{0.55}    % Layer 2 mid-tone
\definecolor{cL2dark}{gray}{0.30}   % Layer 2 darkest
\definecolor{cBoundary}{gray}{0.40}
\definecolor{cGrid}{gray}{0.88}
\definecolor{cText}{gray}{0.15}
\definecolor{cAnnot}{gray}{0.60}
\definecolor{cGapBorder}{gray}{0.65}
\definecolor{cGapFill}{gray}{0.95}
\definecolor{cGapText}{gray}{0.30}

% ═══════════════════════════════════════════════════════════════
% Panel (a): Cascade architecture (full width)
% ═══════════════════════════════════════════════════════════════
\textbf{(a)}\par\vspace{0.3em}
\begin{tikzpicture}[
    x=1cm, y=1cm,
    flowbox/.style={draw=#1!40!black, rounded corners=3pt, fill=#1,
                    minimum width=1.7cm, minimum height=0.65cm,
                    font=\scriptsize\bfseries, text=white,
                    line width=0.3pt, align=center},
    metricbadge/.style={rounded corners=2pt, fill=#1, opacity=0.85,
                        font=\tiny\bfseries, text=white,
                        inner sep=2.5pt},
    arr/.style={-{Stealth[length=1.5mm, width=1mm]}, #1, line width=0.6pt},
]

% ── Layer 1 background band ──
\fill[cL1light, rounded corners=4pt]
    (0, 2.5) rectangle (9.5, 5.0);
\draw[cL1, line width=0.5pt, rounded corners=4pt]
    (0, 2.5) rectangle (9.5, 5.0);
\node[font=\tiny\itshape, text=cL1, anchor=west] at (0.15, 4.8)
    {Layer 1: Sequence-determined};

% Layer 1 boxes
\node[flowbox=cL1]
    (protseq) at (1.2, 4.15) {Protein\\sequence};
\node[flowbox=cL1mid]
    (sitepos) at (1.2, 3.15) {Site\\position};
\node[flowbox=cL1]
    (esm2) at (3.6, 3.65) {ESM-2 +\\retrieval\\model};
\node[flowbox=cL1dark]
    (cluster) at (6.0, 3.65) {Cluster\\(K\,=\,20)};
\node[metricbadge=cL1] at (8.0, 3.65) {H@5 = 0.786};

% Layer 1 arrows
\draw[arr=cL1] (protseq.east) -- (esm2.170);
\draw[arr=cL1] (sitepos.east) -- (esm2.190);
\draw[arr=cL1] (esm2.east) -- (cluster.west);

% ── Layer 2 background band ──
\fill[cL2light, rounded corners=4pt]
    (0, 0.0) rectangle (9.5, 2.2);
\draw[cL2, line width=0.5pt, rounded corners=4pt]
    (0, 0.0) rectangle (9.5, 2.2);
\node[font=\tiny\itshape, text=cL2, anchor=west] at (0.15, 2.0)
    {Layer 2: Environment-determined};

% Layer 2 boxes
\node[flowbox=cL2]
    (celltype) at (1.2, 1.05) {Cell-type\\enzyme expr.};
\node[flowbox=cL2mid]
    (scoring) at (3.6, 1.05) {Enzyme\\scoring};
\node[flowbox=cL2dark]
    (glycanid) at (6.0, 1.05) {Glycan ID};
\node[metricbadge=cL2] at (8.0, 1.05) {MRR = 0.066};

% Layer 2 arrows
\draw[arr=cL2] (celltype.east) -- (scoring.west);
\draw[arr=cL2] (scoring.east)  -- (glycanid.west);

% ── Cross-boundary arrow (dashed) ──
\draw[-{Stealth[length=1.5mm, width=1mm]}, cBoundary,
      line width=0.7pt, dashed]
    (cluster.south) .. controls +(0,-0.5) and +(0.8,0.7) .. (scoring.north);

% ── Data-gap callout ──
\node[draw=cGapBorder, rounded corners=3pt,
      fill=cGapFill, line width=0.4pt,
      font=\tiny, text=cGapText, align=center]
    at (10.5, 1.05) {39\% glycans with\\enzyme annotations};

\end{tikzpicture}

\vspace{1.2em}

% ═══════════════════════════════════════════════════════════════
% Bottom row: (b) left, (c) right
% ═══════════════════════════════════════════════════════════════
\begin{minipage}[t]{0.47\textwidth}
\centering

% ═══════════════════════════════════════════════════════════════
% Panel (b): Two-layer conceptual model
% ═══════════════════════════════════════════════════════════════
\textbf{(b)}\par\vspace{0.3em}
\begin{tikzpicture}[x=0.95cm, y=1cm]

% ── Layer 1 band ──
\fill[cL1light, rounded corners=4pt] (1.2, 2.6) rectangle (7.0, 4.6);
\draw[cL1, line width=0.6pt, rounded corners=4pt] (1.2, 2.6) rectangle (7.0, 4.6);
\node[font=\scriptsize\bfseries, text=cL1] at (4.1, 4.2)
    {Layer 1: Sequence-determined};
\node[font=\tiny, text=cText, align=center] at (4.1, 3.4)
    {Site classification (F1 = 0.924)\\
     Type prediction (Acc = 0.844)\\
     Structural cluster (H@5 = 0.786)};

% ── Boundary dashed line ──
\draw[cBoundary, dashed, line width=0.7pt, opacity=0.7] (1.3, 2.35) -- (6.9, 2.35);

% ── Layer 2 band ──
\fill[cL2light, rounded corners=4pt] (1.2, 0.3) rectangle (7.0, 2.1);
\draw[cL2, line width=0.6pt, rounded corners=4pt] (1.2, 0.3) rectangle (7.0, 2.1);
\node[font=\scriptsize\bfseries, text=cL2] at (4.1, 1.7)
    {Layer 2: Environment-determined};
\node[font=\tiny, text=cText, align=center] at (4.1, 0.95)
    {Exact glycan ID (MRR = 0.066--0.118)\\
     Enzyme expr.\;\;|\;\;Golgi transit\;\;|\;\;NTP sugars};

% ── 10× gap annotation ──
\draw[{Stealth[length=1.5mm]}-{Stealth[length=1.5mm]}, cText, line width=0.8pt]
    (7.8, 1.2) -- (7.8, 3.6);
\node[font=\scriptsize\bfseries, text=cText, fill=white, inner sep=1pt]
    at (7.8, 2.4) {10$\times$\,gap};

% ── Side annotations ──
\node[font=\tiny\itshape, text=cL1, align=center, anchor=east] at (1.0, 3.6)
    {Predictable\\from sequence};
\node[font=\tiny\itshape, text=cL2, align=center, anchor=east] at (1.0, 1.2)
    {Requires\\cellular context};

\end{tikzpicture}
\end{minipage}%
\hfill%
\begin{minipage}[t]{0.50\textwidth}
\centering

% ═══════════════════════════════════════════════════════════════
% Panel (c): Performance across hierarchy — horizontal bars
% ═══════════════════════════════════════════════════════════════
\textbf{(c)}\par\vspace{0.3em}
\begin{tikzpicture}[x=1cm, y=1cm]

\pgfmathsetmacro{\bs}{5.8}

% ── Layer background spans ──
\fill[cL1light, opacity=0.5] (0, 1.4) rectangle (\bs+0.6, 5.3);
\fill[cL2light, opacity=0.5] (0, -0.35) rectangle (\bs+0.6, 1.4);

% ── Grid lines ──
\foreach \gx in {0,0.2,...,1.0}
    \draw[cGrid, line width=0.2pt] (\gx*\bs, -0.3) -- (\gx*\bs, 5.2);

% ── Layer labels ──
\node[font=\tiny\bfseries, text=cL1, anchor=east] at (\bs+0.5, 5.1) {Layer 1};
\node[font=\tiny\bfseries, text=cL2, anchor=east] at (\bs+0.5, -0.15) {Layer 2};

% ── Bars (Layer 1: dark gray) ──
\fill[cL1] (0, 4.55) rectangle (0.924*\bs, 4.95);
\node[anchor=east, font=\tiny, text=cText] at (-0.1, 4.75) {N-linked classif.};
\node[anchor=west, font=\tiny, text=cText] at (0.924*\bs+0.06, 4.75) {F1\,=\,0.924};

\fill[cL1] (0, 3.85) rectangle (0.683*\bs, 4.25);
\node[anchor=east, font=\tiny, text=cText] at (-0.1, 4.05) {Site ranking};
\node[anchor=west, font=\tiny, text=cText] at (0.683*\bs+0.06, 4.05) {R@3\,=\,0.683};

\fill[cL1] (0, 3.15) rectangle (0.844*\bs, 3.55);
\node[anchor=east, font=\tiny, text=cText] at (-0.1, 3.35) {Glycosylation type};
\node[anchor=west, font=\tiny, text=cText] at (0.844*\bs+0.06, 3.35) {Acc\,=\,0.844};

\fill[cL1] (0, 2.45) rectangle (0.858*\bs, 2.85);
\node[anchor=east, font=\tiny, text=cText] at (-0.1, 2.65) {Structural family};
\node[anchor=west, font=\tiny, text=cText] at (0.858*\bs+0.06, 2.65) {Top-2\,=\,0.858};

\fill[cL1] (0, 1.75) rectangle (0.786*\bs, 2.15);
\node[anchor=east, font=\tiny, text=cText] at (-0.1, 1.95) {Structural cluster};
\node[anchor=west, font=\tiny, text=cText] at (0.786*\bs+0.06, 1.95) {H@5\,=\,0.786};

% ── Bars (Layer 2: medium gray) ──
\fill[cL2] (0, 0.7) rectangle (0.066*\bs, 1.1);
\node[anchor=east, font=\tiny, text=cText] at (-0.1, 0.9) {Exact ID (site)};
\node[anchor=west, font=\tiny, text=cText] at (0.066*\bs+0.06, 0.9) {MRR\,=\,0.066};

\fill[cL2] (0, 0.0) rectangle (0.118*\bs, 0.4);
\node[anchor=east, font=\tiny, text=cText] at (-0.1, 0.2) {Exact ID (prot.)};
\node[anchor=west, font=\tiny, text=cText] at (0.118*\bs+0.06, 0.2) {MRR\,=\,0.118};

% ── X-axis ──
\draw[cText, line width=0.3pt] (0,-0.35) -- (\bs,-0.35);
\foreach \tick/\lab in {0/0, 0.2/0.2, 0.4/0.4, 0.6/0.6, 0.8/0.8, 1.0/1.0}
    \draw[cText, line width=0.2pt] (\tick*\bs, -0.35) -- (\tick*\bs, -0.42)
    node[below, font=\tiny, text=cText] {\lab};
\node[font=\tiny, text=cText] at (\bs*0.5, -0.8) {Performance};

\end{tikzpicture}
\end{minipage}

\caption{\textbf{Two-layer determination model and cell-type analysis.}
(a)~Cascade architecture combining sequence-based cluster prediction (Layer~1)
with cell-type enzyme expression scoring (Layer~2).
(b)~Conceptual model of glycosylation determination: protein sequence determines
structural class (Layer~1; H@5 = 0.786), while exact glycan identity requires
cellular context (Layer~2; MRR = 0.066--0.118). The 10-fold gap persists across
model architectures.
(c)~Performance across the prediction hierarchy. Dark bars represent
sequence-determined tasks (Layer~1); medium-gray bars represent
environment-dependent tasks (Layer~2).}
\label{fig:twolayer}
\end{figure}


% ══════════════════════════════════════════════════════════════════
% References
% ══════════════════════════════════════════════════════════════════

\begin{thebibliography}{27}

\bibitem{varki2017}
Varki,A. (2017) Biological roles of glycans. \textit{Glycobiology}, \textbf{27}, 3--49.

\bibitem{reily2019}
Reily,C. \textit{et al.} (2019) Glycosylation in health and disease. \textit{Nat. Rev. Nephrol.}, \textbf{15}, 346--366.

\bibitem{schjoldager2020}
Schjoldager,K.T. \textit{et al.} (2020) Global view of human protein glycosylation pathways and functions. \textit{Nat. Rev. Mol. Cell Biol.}, \textbf{21}, 729--749.

\bibitem{narimatsu2019}
Narimatsu,Y. \textit{et al.} (2019) An atlas of human glycosylation pathways enables display of the human glycome by gene engineering. \textit{Mol. Cell}, \textbf{75}, 394--407.

\bibitem{stanley2022}
Stanley,P. \textit{et al.} (2022) N-Glycans. In: Varki,A. \textit{et al.} (eds) \textit{Essentials of Glycobiology}, 4th edn, Ch.~9. Cold Spring Harbor Laboratory Press.

\bibitem{varki2022essentials}
Varki,A. \textit{et al.} (2022) \textit{Essentials of Glycobiology}, 4th edn. Cold Spring Harbor Laboratory Press.

\bibitem{tiemeyer2017}
Tiemeyer,M. \textit{et al.} (2017) GlyTouCan: an accessible glycan structure repository. \textit{Glycobiology}, \textbf{27}, 915--919.

\bibitem{uniprot2023}
The UniProt Consortium (2023) UniProt: the Universal Protein Knowledgebase in 2023. \textit{Nucleic Acids Res.}, \textbf{51}, D523--D531.

\bibitem{york2020}
York,W.S. \textit{et al.} (2020) GlyGen: computational and informatics resources for glycoscience. \textit{Glycobiology}, \textbf{30}, 72--73.

\bibitem{zdrazil2024}
Zdrazil,B. \textit{et al.} (2024) The ChEMBL Database in 2023. \textit{Nucleic Acids Res.}, \textbf{52}, D1180--D1192.

\bibitem{alocci2019}
Alocci,D. \textit{et al.} (2019) GlyConnect: glycoproteomics goes visual, interactive, and analytical. \textit{J. Proteome Res.}, \textbf{18}, 698--708.

\bibitem{burkholz2022}
Burkholz,R. \textit{et al.} (2022) SweetNet: a neural network for predicting glycan properties from graph representation. \textit{Bioinformatics}, \textbf{38}, 4599--4605.

\bibitem{thomes2024}
Thom\`es,L. and Burkholz,R. (2024) GIFFLAR: a higher-order graph neural network for glycan function prediction. \textit{Preprint at bioRxiv}.

\bibitem{thomes2021}
Thom\`es,L. \textit{et al.} (2021) glycowork: a Python package for glycan data science and machine learning. \textit{Glycobiology}, \textbf{31}, 1240--1244.

\bibitem{lin2023}
Lin,Z. \textit{et al.} (2023) Evolutionary-scale prediction of atomic-level protein structure with a language model. \textit{Science}, \textbf{379}, 1123--1130.

\bibitem{bordes2013}
Bordes,A. \textit{et al.} (2013) Translating embeddings for modeling multi-relational data. \textit{Adv. Neural Inf. Process. Syst.}, \textbf{26}.

\bibitem{tabulasapiens2022}
The Tabula Sapiens Consortium (2022) The Tabula Sapiens: a multiple-organ, single-cell transcriptomic atlas of humans. \textit{Science}, \textbf{376}, eabl4896.

\bibitem{uhlen2015}
Uhl\'en,M. \textit{et al.} (2015) Tissue-based map of the human proteome. \textit{Science}, \textbf{347}, 1260419.

\bibitem{minguez2015}
Minguez,P. \textit{et al.} (2015) PTMcode v2: a resource for functional associations of post-translational modifications within and between proteins. \textit{Nucleic Acids Res.}, \textbf{43}, D503--D509.

\bibitem{kellman2024}
Kellman,B.P. \textit{et al.} (2024) Elucidating Human Protein Glycosylation using an Informed Search of Artificial Neural Network Ensembles (InSaNNE). \textit{Preprint at bioRxiv}, doi:10.1101/2024.02.02.578654.

\bibitem{bojar2022}
Bojar,D. and Lisacek,F. (2022) Glycoinformatics in the Artificial Intelligence Era. \textit{Chem.\ Rev.}, \textbf{122}, 15971--15988.

\bibitem{congly2025}
Li,M. \textit{et al.} (2025) CoNglyPred: accurate prediction of protein N-linked glycosylation sites combining ESM-2 and AlphaFold2 features. \textit{PROTEOMICS}, \textbf{25}, e2400129.

\bibitem{glyml2024}
Bojar,D. \textit{et al.} (2024) GlycanML: A Multi-Task and Multi-Structure Benchmark for Glycan Machine Learning. \textit{Preprint at arXiv}, arXiv:2405.16206.

\bibitem{yadav2022}
Yadav,V. \textit{et al.} (2022) Golgi glycan processing as an optimal information coding channel. \textit{eLife}, \textbf{11}, e76505.

\bibitem{chrysinas2024}
Chrysinas,P. \textit{et al.} (2024) Single cell analysis of human glycosyltransferases and glycosidases. \textit{NAR Genomics Bioinformatics}, \textbf{6}, lqae069.

\bibitem{glycoimpact2024}
Lundstr\"om,J. \textit{et al.} (2024) GlycoImpact: In-silico glycan scoring by amino acid substitution matrix. \textit{Preprint at bioRxiv}, doi:10.1101/2024.09.16.613237.

\end{thebibliography}

\end{document}
